
\medskip
\phantomsection\label{prop:sga3-VIB-7.4}%
\noindent\textbf{Proposition 7.4.---}
\emph{Let \(G\) be a \(k\)-group locally of finite type, let
\(X\) be a finite-type \(k\)-scheme that is geometrically
reduced and geometrically connected, and let \(f:X\to G\) be a
\(k\)-morphism such that \(f(X)\) contains the unit element
\(e\) of \(G\). Then, with the notation of
\hyperref[concl:sga3-VIB-7.1-ii]{7.1\textup{(ii)}}, there
exists an integer \(N\) such that, set-theoretically,}
\phantomsection\label{form:sga3-VIB-7.4-stabilization}%
\[
  f^N(X^N)=\Gamma_G(f).
\]

\emph{Proof.}
By \hyperref[concl:sga3-VIB-7.1-iii]{7.1\textup{(iii)}} and
EGA~IV\(_2\), 2.6.1, we may assume that \(k\) is algebraically
closed. By
\hyperref[cor:sga3-VIB-7.2.1]{Corollary 7.2.1}, we may assume
that \(G=G^0\). Finally, with the notation of
\hyperref[rem:sga3-VIB-7.2-v]{7.2\textup{(v)}}, it is enough to
show that there exists an integer \(N\) such that
\[
  f'^{\,N}(X'^{\,N})=\Gamma_G(f).
\]

\emph{First case.}
Suppose that \(X\) is irreducible. The closed irreducible
subsets
\[
  \overline{f'^{\,n}(X'^{\,n})}
\]
then form an increasing sequence in the space \(G\), which is
Noetherian because \(G=G^0\) is of finite type over \(k\)
(\hyperref[prop:sga3-VIA-2.4]{VI\(_A\), 2.4}). Hence this
sequence is stationary, and there exists an integer \(m\) such
that
\[
  \overline{f'^{\,m}(X'^{\,m})}=\Gamma_G(f).
\]

Moreover, because \(X\) and \(G\) are of finite type over
\(k\), the morphisms \(f'^{\,n}\) are of finite type over
\(k\). Consequently, \(f'^{\,m}(X'^{\,m})\) is constructible
in \(G\) (EGA~IV\(_1\), 1.8.5), and therefore contains an open
subset \(U\) of its closure \(\Gamma_G(f)\)
(EGA~0\(_{\mathrm{IV}}\), 9.2.3). By
\hyperref[sga3:VIA:0.5]{VI\(_A\), 0.5}, one then has
\phantomsection\label{form:sga3-VIB-7.4-first-case-chain}%
\[
  \Gamma_G(f)
  \subset U\cdot U
  \subset f'^{\,2m}(X'^{\,2m})
  \subset \Gamma_G(f),
\]
and hence
\[
  f'^{\,2m}(X'^{\,2m})=\Gamma_G(f).
\]

\emph{Second case.}
Suppose that \(X\) has exactly two irreducible components
\(A_1\) and \(A_2\). Since \(X\) is connected and \(k\) is
algebraically closed, there exists
\[
  a\in(A_1\cap A_2)(k).
\]
The four irreducible subsets
\[
  A_i\times_kA_j
  \qquad (i,j=1,2)
\]
therefore cover \(X\times_kX\), and the image of each under the
morphism
\[
  f'^{\,1}=\mu\circ\bigl(f\times_k(c\circ f)\bigr)
\]
contains \(e\). Let \(f'^{\,1}_{ij}\) denote the restriction of
\(f'^{\,1}\) to the reduced subscheme \(A_i\times_kA_j\), and
put
\[
\begin{aligned}
  Y={}&
  (A_1\times_kA_1)
  \times_k(A_1\times_kA_2)
  \times_k(A_2\times_kA_2)
  \times_k(A_2\times_kA_1),\\
  g={}&
  \mu\circ\Bigl(
    \mu\circ
      (f'^{\,1}_{11}\times_k f'^{\,1}_{12})
    \times_k
    \mu\circ
      (f'^{\,1}_{22}\times_k f'^{\,1}_{21})
  \Bigr).
\end{aligned}
\]
The scheme \(Y\) is irreducible, reduced, and of finite type.
By the first case, there is therefore an integer \(m\) such
that
\[
  g'^{\,m}(Y'^{\,m})=\Gamma_G(g).
\]
For every \(n\geq1\), however,\begingroup
\renewcommand{\thefootnote}{\fnsymbol{footnote}}%
\footnote[1]{\label{note:sga3-VIB-7.4-Yprime}Translator's note:
the French source prints \(g'^{\,n}(Y^{\,n})\) in the middle
term. Since the source of \(g'^{\,n}\) is \(Y'^{\,n}\), the
prime is restored here; compare the preceding equality
\(g'^{\,m}(Y'^{\,m})=\Gamma_G(g)\).}%
\endgroup
\phantomsection\label{form:sga3-VIB-7.4-two-component-chain}%
\[
  f'^{\,n}(X'^{\,n})
  \subset g'^{\,n}(Y'^{\,n})
  \subset f'^{\,4n}(X'^{\,4n}).
\]
It follows that
\[
  \Gamma_G(f)=\Gamma_G(g)
  \qquad\text{and}\qquad
  f'^{\,4m}(X'^{\,4m})=\Gamma_G(f).
\]

\emph{General case.}
We argue by induction on the number of irreducible components
of \(X\); this number is finite because \(X\), being of finite
type over \(k\), is Noetherian. Suppose the proposition has
been proved when \(X\) has at most \(r-1\) irreducible
components, and suppose that \(X\) has \(r\) such components,
say \(A_1,\ldots,A_r\). The unit \(e\) belongs to the image of
one of the \(A_i\); suppose, for example, that
\[
  e\in f(A_1).
\]
Put
\[
  X_r=A_1\cup\cdots\cup A_{r-1},
  \qquad
  Y=\Gamma_G(f_r),
\]
where \(f_r\) denotes the restriction of \(f\) to the reduced
subscheme \(X_r\). We choose the numbering of the \(A_i\) so
that \(X_r\) is connected, which is always possible. Then
\(Y\) is a closed, reduced, irreducible subgroup of \(G\), by
\hyperref[cor:sga3-VIB-7.2.1]{Corollary 7.2.1}.

Put
\[
  Z=\overline{f(A_r)}
\]
and let \(T=Y\cup Z\), endowed with its closed reduced
subscheme structure; let \(i:T\hookrightarrow G\) be the
inclusion. It follows immediately from
\hyperref[rem:sga3-VIB-7.2-ii]{7.2\textup{(ii)}} that
\[
  \Gamma_G(i)=\Gamma_G(f).
\]
Moreover, \(T\) is connected: since \(X\) is connected,
\(X_r\) and \(A_r\) have a common point \(a\), and \(Y\) and
\(Z\) have the common point \(f(a)\). Also, \(e\in T\), and
\(T\) has at most two irreducible components, because \(Y\)
and \(Z\) are irreducible.

By the induction hypothesis, there exists an integer \(m\)
such that
\[
  f_r'^{\,m}(X_r'^{\,m})
  =\Gamma_G(f_r)=Y.
\]
Since \(X=X_r\cup A_r\), and since
\[
  Z=\overline{f(A_r)}
  \subset\overline{f'^{\,m}(X'^{\,m})}
\]
because \(e\in f(A_1)\), we obtain
\phantomsection\label{form:sga3-VIB-7.4-general-case-inclusion}%
\[
  f(X)
  \subset Y\cup Z
  \subset\overline{f'^{\,m}(X'^{\,m})}.
\]

On the other hand, because \(T\) has at most two irreducible
components, the second case gives an integer \(p\) such that
\[
  i'^{\,p}(T'^{\,p})
  =\Gamma_G(i)=\Gamma_G(f).
\]
Now
\[
  T\subset\overline{f'^{\,m}(X'^{\,m})},
\]
and hence
\[
  \overline{f'^{\,mp}(X'^{\,mp})}=\Gamma_G(f).
\]
As in the first case, this last equality implies
\[
  f'^{\,2mp}(X'^{\,2mp})=\Gamma_G(f).
\]
This completes the proof.
